\section{Allocated vs. Unallocated Space}

\begin{description}
	\item[Allocated Space] Aktiv genutzter Speicher mit Referenz im Dateisystem.
		Dateien sind über das Dateisystem direkt zugreifbar.
	\item[Unallocated Space] Freigegebener Speicher — Daten physisch noch vorhanden, bis sie überschrieben werden.
		Forensisch besonders wertvoll, da gelöschte Dateien hier oft noch rekonstruierbar sind.
\end{description}

\section{EWF-Format (E01)}

Expert Witness Format — das Standard-Imageformat in der digitalen Forensik.

\begin{itemize}
	\item Komprimiertes forensisches Image mit integrierten Metadaten (Case-Nr., Examiner, Datum)
	\item Integrierte Prüfsummen (MD5) zur Integritätssicherung
	\item Enthält sowohl Allocated als auch Unallocated Space (bitgenaue Kopie)
	\item Segmente werden bei 2 GB gesplittet (\texttt{.E01}–\texttt{.E07})
	\item Metadaten einsehen mit \texttt{ewfinfo}, Integrität prüfen mit \texttt{ewfverify}
\end{itemize}

\section{RAW → EWF → xmount Kreislauf}

\begin{enumerate}
	\item Originaldisk (z.B. 8 TB) wird mit \texttt{ewfacquire} als komprimiertes EWF-Image gesichert (~4 TB)
	\item \texttt{xmount} stellt das EWF als virtuelles RAW-Image (8 TB) bereit
	\item Ein Overlay (\texttt{--cache}) speichert alle Änderungen — das Original bleibt unverändert
\end{enumerate}

\begin{lstlisting}[language=bash]
ewfacquire /dev/sdX              # Disk -> EWF
xmount --in ewf img.E0* --cache /tmp/img.ovl --out raw /ewf
# /ewf/img.dd ist jetzt ein virtuelles RAW-Image
\end{lstlisting}

\section{Physische vs. logische Analyse}

\begin{description}
	\item[Physische Analyse] Zugriff auf alles (Allocated + Unallocated).
		Über TSK-Tools wie \texttt{fls}, \texttt{icat}, \texttt{istat}.
		Ermöglicht Zugriff auf gelöschte Dateien und Dateifragmente.
	\item[Logische Analyse] Image wird als Dateisystem gemountet (\texttt{mount -o ro}).
		Nur Allocated Space sichtbar — wie ein normales Dateisystem.
		Einfacher zu navigieren, aber weniger Informationen.
\end{description}

\section{File Carving}

Letzte Möglichkeit zur Datenrettung — Dateien werden anhand bekannter Header/Footer-Signaturen rekonstruiert, ohne Dateisystem-Metadaten.

\begin{itemize}
	\item Tools: \texttt{foremost}, \texttt{scalpel}
	\item Verliert Dateinamen, Verzeichnisstruktur und Metadaten
	\item Funktioniert auch bei beschädigten oder formatierten Dateisystemen
	\item Nur sinnvoll, wenn \texttt{tsk\_recover} keine Ergebnisse liefert
\end{itemize}

\section{Chain of Custody \& Write Blocker}

\begin{description}
	\item[Chain of Custody] Lückenlose Dokumentation aller Aktionen an Beweismitteln.
		Jede Analyse, jeder Zugriff muss protokolliert werden, um die Beweiskraft zu erhalten.
	\item[Write Blocker] Hardware oder Software, die Schreibzugriffe auf Beweismittel verhindert.
		EWF-Images sind inhärent read-only.
		Bei \texttt{xmount} sorgt der Overlay-Cache dafür, dass das Original unverändert bleibt.
\end{description}

\section{Integrität}

Immer als ersten Schritt die Integrität prüfen — vor und nach der Analyse.

\begin{lstlisting}[language=bash]
ewfverify image.E01              # Interne EWF-Pruefsummen pruefen
md5sum image.E01 image.E02       # MD5-Hashes berechnen
md5sum -c md5s.txt               # Gegen gespeicherte Hashes vergleichen
\end{lstlisting}
